% Options for packages loaded elsewhere
\PassOptionsToPackage{unicode}{hyperref}
\PassOptionsToPackage{hyphens}{url}
\PassOptionsToPackage{dvipsnames,svgnames*,x11names*}{xcolor}
%
\documentclass[
]{article}
\usepackage{amsmath,amssymb}
\usepackage{lmodern}
\usepackage{iftex}
\ifPDFTeX
  \usepackage[T1]{fontenc}
  \usepackage[utf8]{inputenc}
  \usepackage{textcomp} % provide euro and other symbols
\else % if luatex or xetex
  \usepackage{unicode-math}
  \defaultfontfeatures{Scale=MatchLowercase}
  \defaultfontfeatures[\rmfamily]{Ligatures=TeX,Scale=1}
\fi
% Use upquote if available, for straight quotes in verbatim environments
\IfFileExists{upquote.sty}{\usepackage{upquote}}{}
\IfFileExists{microtype.sty}{% use microtype if available
  \usepackage[]{microtype}
  \UseMicrotypeSet[protrusion]{basicmath} % disable protrusion for tt fonts
}{}
\makeatletter
\@ifundefined{KOMAClassName}{% if non-KOMA class
  \IfFileExists{parskip.sty}{%
    \usepackage{parskip}
  }{% else
    \setlength{\parindent}{0pt}
    \setlength{\parskip}{6pt plus 2pt minus 1pt}}
}{% if KOMA class
  \KOMAoptions{parskip=half}}
\makeatother
\usepackage{xcolor}
\IfFileExists{xurl.sty}{\usepackage{xurl}}{} % add URL line breaks if available
\IfFileExists{bookmark.sty}{\usepackage{bookmark}}{\usepackage{hyperref}}
\hypersetup{
  colorlinks=true,
  linkcolor={blue},
  filecolor={Maroon},
  citecolor={Blue},
  urlcolor={blue},
  pdfcreator={LaTeX via pandoc}}
\urlstyle{same} % disable monospaced font for URLs
\usepackage[margin=1in]{geometry}
\usepackage{color}
\usepackage{fancyvrb}
\newcommand{\VerbBar}{|}
\newcommand{\VERB}{\Verb[commandchars=\\\{\}]}
\DefineVerbatimEnvironment{Highlighting}{Verbatim}{commandchars=\\\{\}}
% Add ',fontsize=\small' for more characters per line
\newenvironment{Shaded}{}{}
\newcommand{\AlertTok}[1]{\textcolor[rgb]{1.00,0.00,0.00}{\textbf{#1}}}
\newcommand{\AnnotationTok}[1]{\textcolor[rgb]{0.38,0.63,0.69}{\textbf{\textit{#1}}}}
\newcommand{\AttributeTok}[1]{\textcolor[rgb]{0.49,0.56,0.16}{#1}}
\newcommand{\BaseNTok}[1]{\textcolor[rgb]{0.25,0.63,0.44}{#1}}
\newcommand{\BuiltInTok}[1]{#1}
\newcommand{\CharTok}[1]{\textcolor[rgb]{0.25,0.44,0.63}{#1}}
\newcommand{\CommentTok}[1]{\textcolor[rgb]{0.38,0.63,0.69}{\textit{#1}}}
\newcommand{\CommentVarTok}[1]{\textcolor[rgb]{0.38,0.63,0.69}{\textbf{\textit{#1}}}}
\newcommand{\ConstantTok}[1]{\textcolor[rgb]{0.53,0.00,0.00}{#1}}
\newcommand{\ControlFlowTok}[1]{\textcolor[rgb]{0.00,0.44,0.13}{\textbf{#1}}}
\newcommand{\DataTypeTok}[1]{\textcolor[rgb]{0.56,0.13,0.00}{#1}}
\newcommand{\DecValTok}[1]{\textcolor[rgb]{0.25,0.63,0.44}{#1}}
\newcommand{\DocumentationTok}[1]{\textcolor[rgb]{0.73,0.13,0.13}{\textit{#1}}}
\newcommand{\ErrorTok}[1]{\textcolor[rgb]{1.00,0.00,0.00}{\textbf{#1}}}
\newcommand{\ExtensionTok}[1]{#1}
\newcommand{\FloatTok}[1]{\textcolor[rgb]{0.25,0.63,0.44}{#1}}
\newcommand{\FunctionTok}[1]{\textcolor[rgb]{0.02,0.16,0.49}{#1}}
\newcommand{\ImportTok}[1]{#1}
\newcommand{\InformationTok}[1]{\textcolor[rgb]{0.38,0.63,0.69}{\textbf{\textit{#1}}}}
\newcommand{\KeywordTok}[1]{\textcolor[rgb]{0.00,0.44,0.13}{\textbf{#1}}}
\newcommand{\NormalTok}[1]{#1}
\newcommand{\OperatorTok}[1]{\textcolor[rgb]{0.40,0.40,0.40}{#1}}
\newcommand{\OtherTok}[1]{\textcolor[rgb]{0.00,0.44,0.13}{#1}}
\newcommand{\PreprocessorTok}[1]{\textcolor[rgb]{0.74,0.48,0.00}{#1}}
\newcommand{\RegionMarkerTok}[1]{#1}
\newcommand{\SpecialCharTok}[1]{\textcolor[rgb]{0.25,0.44,0.63}{#1}}
\newcommand{\SpecialStringTok}[1]{\textcolor[rgb]{0.73,0.40,0.53}{#1}}
\newcommand{\StringTok}[1]{\textcolor[rgb]{0.25,0.44,0.63}{#1}}
\newcommand{\VariableTok}[1]{\textcolor[rgb]{0.10,0.09,0.49}{#1}}
\newcommand{\VerbatimStringTok}[1]{\textcolor[rgb]{0.25,0.44,0.63}{#1}}
\newcommand{\WarningTok}[1]{\textcolor[rgb]{0.38,0.63,0.69}{\textbf{\textit{#1}}}}
\usepackage{longtable,booktabs,array}
\usepackage{calc} % for calculating minipage widths
% Correct order of tables after \paragraph or \subparagraph
\usepackage{etoolbox}
\makeatletter
\patchcmd\longtable{\par}{\if@noskipsec\mbox{}\fi\par}{}{}
\makeatother
% Allow footnotes in longtable head/foot
\IfFileExists{footnotehyper.sty}{\usepackage{footnotehyper}}{\usepackage{footnote}}
\makesavenoteenv{longtable}
\setlength{\emergencystretch}{3em} % prevent overfull lines
\providecommand{\tightlist}{%
  \setlength{\itemsep}{0pt}\setlength{\parskip}{0pt}}
\setcounter{secnumdepth}{5}
\ifLuaTeX
  \usepackage{selnolig}  % disable illegal ligatures
\fi

\author{}
\date{}

\begin{document}

{
\hypersetup{linkcolor=}
\setcounter{tocdepth}{3}
\tableofcontents
}
\hypertarget{varblockspmm-a-shape-aware-cuda-kernel-for-packed-variable-block-sparsedense-matrix-multiplication}{%
\section{VarBlockSpMM: A Shape-Aware CUDA Kernel for Packed
Variable-Block Sparse--Dense Matrix
Multiplication}\label{varblockspmm-a-shape-aware-cuda-kernel-for-packed-variable-block-sparsedense-matrix-multiplication}}

\textbf{Cohen}\\
\emph{Project research manuscript, 21 August 2026}

\hypertarget{abstract}{%
\subsection{Abstract}\label{abstract}}

Sparse--dense matrix multiplication (SpMM) is commonly implemented using
scalar sparse formats or fixed-size block formats. Both choices are
imperfect when a matrix naturally consists of dense blocks whose row and
column dimensions vary independently: scalar expansion discards block
structure, while global padding introduces storage and bandwidth
overhead. This paper presents \textbf{VarBlockSpMM}, an FP32 CUDA
implementation of the non-transpose product \(C=AB\) for a packed
variable-block sparse row (VBSR) matrix \(A\) and dense column-major
panels \(B\) and \(C\). Block heights and widths independently take
values in \(\{8,16,\ldots,64\}\), and no global block-size padding is
required.

The implementation combines row-owned output assignment,
right-hand-side-width specialization, per-thread instruction-level
parallelism, cooperative staging of reusable \(B\) tiles, and a measured
row-shape dispatch. For 32-column panels, tall rows use asynchronous
double buffering while short rows retain a lower-overhead single-buffer
kernel. The direct path requires no atomics or execute-time temporary
workspace. It is evaluated against a persistent scalar-CSR cuSPARSE
baseline and a persistent slot-grouped cuBLAS baseline over 128
synthetic regimes on one NVIDIA GeForce RTX 5060 Ti. In the checked-in
short sweep, the final direct method is fastest in all 128 cases.
Relative to the earlier four-accumulator direct kernel, successive
wider-ILP and staging changes provide a 1.46× geometric-mean speedup
over the complete grid; the final RHS-32 double-buffer refinement adds
1.11× over its 32-case slice. Profiling links the improvement to reduced
redundant loads and reduced long-scoreboard stalls. These results
establish a strong single-device proof of concept, but not a general
performance claim: multi-architecture testing, real application
matrices, repeated-seed experiments, and stronger statistical reporting
remain necessary.

\hypertarget{introduction}{%
\subsection{1. Introduction}\label{introduction}}

Many hierarchical, finite-element, graph, and scientific operators
contain small dense submatrices embedded in a sparse global structure.
When these submatrices share a single size, block sparse row (BSR)
storage exposes useful locality. Real partitions, however, often yield
blocks with different heights and widths. Converting such a matrix to
scalar CSR loses explicit block structure; rounding every block to one
maximum dimension wastes storage and arithmetic; and dispatching a
separate dense GEMM for every block pays substantial scheduling
overhead.

VarBlockSpMM studies the narrower problem

\[
C = A B,
\]

where \(A\in\mathbb{R}^{M\times K}\) is divided into variable-height
block rows and variable-width block columns, while
\(B\in\mathbb{R}^{K\times N}\) and \(C\in\mathbb{R}^{M\times N}\) are
dense column-major panels. The implementation supports FP32,
non-transpose multiplication and \(N\in\{8,16,32,64\}\). Its design
question is: \textbf{can a direct, workspace-free CUDA kernel retain the
compactness of variable blocks while outperforming reasonable persistent
library compositions across a controlled range of irregularity, degree,
locality, and panel width?}

This work makes four engineering contributions:

\begin{enumerate}
\def\labelenumi{\arabic{enumi}.}
\tightlist
\item
  A validated packed VBSR representation with independent row and column
  partitions and 64-bit scalar/value offsets.
\item
  A race-free row-owned kernel family that specializes only panel width
  while retaining variable block dimensions as runtime loop bounds.
\item
  A profiler-guided sequence of optimizations: wide accumulation to
  reuse each \(A\) value, cooperative shared-memory staging of \(B\),
  and selective asynchronous double buffering for tall RHS-32 rows.
\item
  A reproducible comparison against persistent cuSPARSE and
  grouped-cuBLAS constructions, accompanied by correctness, sanitizer,
  timing, and hardware-counter evidence.
\end{enumerate}

The contribution should be interpreted as an implementation and
empirical case study. The associated prior-art search found no directly
matching packed variable-row/variable-column FP32 SpMM interface, but it
was scoped rather than exhaustive and is not a patentability or
universal novelty claim.

\hypertarget{background-and-related-work}{%
\subsection{2. Background and Related
Work}\label{background-and-related-work}}

GPU SpMM research has long emphasized that workload structure and row
irregularity determine the useful mapping of sparse work to warps and
thread blocks. Yang, Buluç, and Owens derive general GPU SpMM design
principles and demonstrate row-splitting and merge-based approaches for
scalar sparse matrices {[}1{]}. More recent systems exploit block
structure and tensor cores, but generally target fixed or structured
block layouts rather than the independently variable blocks studied here
{[}2{]}.

NVIDIA cuSPARSE provides SpMM for several sparse formats. Its
Blocked-ELL path uses a block size belonging to the format and is tuned
for fixed power-of-two blocks {[}3{]}; it does not directly encode a
matrix in which every block row and column may have a different extent.
NVIDIA cuBLAS provides grouped batched GEMM with per-group matrix
dimensions {[}4{]}. This is a closer building block, but independent
GEMMs that update the same output region cannot simply execute together
without respecting accumulation dependencies. VarBlockSpMM therefore
uses a slot-split organization as its cuBLAS baseline: blocks at the
same ordinal position in different block rows form non-overlapping
grouped calls, and slots execute sequentially with accumulation into
\(C\).

The immediate motivating application is adaptive construction of
hierarchical \(H^2\) matrices. Boukaram et al.~describe a non-uniform
batched BSR product and split it into a bounded number of non-uniform
batched GEMM kernels in the absence of a directly suitable GPU
implementation {[}5{]}. VarBlockSpMM targets the underlying packed
variable-block operation rather than the full hierarchical construction
algorithm.

Asynchronous global-to-shared copies are a natural mechanism for hiding
the latency of reusable dense-panel tiles. NVIDIA documents that these
operations decouple transfer initiation from completion and can overlap
data movement with computation; hardware support is available from
compute capability 8.0 for suitably aligned small copies {[}6{]}. The
final RHS-32 tall-row kernel applies this mechanism conservatively,
after profiling showed that a synchronous staged design had shifted its
principal cost from long-scoreboard stalls toward block barriers.

\hypertarget{data-representation-and-semantics}{%
\subsection{3. Data Representation and
Semantics}\label{data-representation-and-semantics}}

Let block-row boundaries be \(r_0=0<r_1<\cdots<r_R=M\) and block-column
boundaries be \(c_0=0<c_1<\cdots<c_S=K\). Block row \(i\) has height
\(h_i=r_{i+1}-r_i\); block column \(j\) has width \(w_j=c_{j+1}-c_j\).
Every stored block \(A_{ij}\) is a dense \(h_i\times w_j\) column-major
array. Its contribution is

\[
C[r_i:r_{i+1},:] \mathrel{+}= A_{ij} B[c_j:c_{j+1},:].
\]

The host format stores block-row pointers, column-block indices, scalar
row/column offsets, value offsets, and tightly packed values. Validation
checks monotonic offsets, legal dimensions, in-range indices, and
agreement between metadata and the value buffer. Device ownership is
separated from execution planning: \texttt{Matrix} owns immutable GPU
storage, while \texttt{Plan} contains a non-owning view and immutable
row-classification metadata. Consequently, the matrix must outlive its
plan.

Packing each block at its actual dimensions avoids the capacity cost of
padding all blocks to \(64\times64\). The tradeoff is that address
generation and loop trip counts remain dynamic inside the kernel. The
format uses 64-bit scalar and value offsets to avoid limiting aggregate
matrix or packed-value size through 32-bit addressing metadata.

\hypertarget{kernel-design}{%
\subsection{4. Kernel Design}\label{kernel-design}}

\hypertarget{row-ownership}{%
\subsubsection{4.1 Row ownership}\label{row-ownership}}

One cooperative thread array (CTA) owns one block row, and its threads
cover disjoint elements of that row's output panel. Because no two CTAs
update the same output scalar, the kernel needs neither atomic
operations nor partial-output reduction storage. Each CTA iterates
through the blocks of its row and accumulates their contributions before
writing its owned elements.

The dispatch specializes on \(N\), not on the runtime block dimensions.
RHS 8 uses a scalar mapping with one output element per thread. For RHS
16 and 32, a thread accumulates eight output columns for one local row;
RHS 64 uses sixteen. Reusing one loaded \(A\) value across multiple
independent accumulators reduces repeated sparse-value traffic and
exposes independent fused multiply-add instructions, at the cost of
registers and reduced thread-level parallelism. Measurements rejected
wider mappings where that tradeoff became unfavorable.

\hypertarget{cooperative-staging}{%
\subsubsection{4.2 Cooperative staging}\label{cooperative-staging}}

For RHS 32 and 64, all local output rows in a CTA consume the same
\(w_j\times N\) slice of \(B\). Direct global loading redundantly
requests these values for each local row. The staged kernels
cooperatively load the active \(B\) slice into shared memory once per
block, then reuse it across local rows. A shared leading dimension of 65
floats separates neighboring columns and avoids the worst power-of-two
bank-alignment pattern. The largest logical RHS-64 tile contains
\(64\times64\) floats, or 16 KiB.

The compiled synchronous staged kernels use 9,344 bytes of shared memory
and 56 registers per thread for RHS 32, and 17,664 bytes and 64
registers per thread for RHS 64, without local-memory spills. RHS 16
remains on direct global loads because its smaller reuse opportunity did
not repay synchronization overhead.

\hypertarget{shape-aware-asynchronous-pipeline}{%
\subsubsection{4.3 Shape-aware asynchronous
pipeline}\label{shape-aware-asynchronous-pipeline}}

RHS-32 rows are classified once during matrix construction. Rows of at
most 16 scalars use a 128-thread single-buffer CTA. Taller rows use 256
threads and two shared tiles. While the CTA computes on tile \(t\),
per-thread asynchronous copies initiate loading of tile \(t+1\). A
pipeline wait and CTA barrier at the block boundary establish visibility
and prevent buffer reuse before computation finishes.

Splitting a mixed-height matrix into two kernel launches is enabled only
when its mean block degree is at least eight. At lower degree,
measurements showed that the second launch costs more than shape
separation saves; homogeneous matrices already require only their
applicable launch. The final tall-row RHS-32 kernel uses 48 registers
per thread and 17,664 bytes of shared memory; its 128-thread companion
uses 54 registers and 9,344 bytes.

The same policy is deliberately not applied to RHS 64. Two full RHS-64
tiles require roughly 34 KiB per CTA. In the representative 4,096-row
case, that allocation reduced residency and regressed median time from
2.891 ms to 4.419 ms. The accepted design therefore retains the
single-buffer RHS-64 kernel.

\hypertarget{baselines-and-experimental-method}{%
\subsection{5. Baselines and Experimental
Method}\label{baselines-and-experimental-method}}

Two persistent baselines separate the proposed representation from
readily available library strategies.

\textbf{Scalar CSR cuSPARSE.} Dense blocks are expanded once into scalar
CSR. The plan owns its CSR arrays, descriptors, library handle, and
workspace, all of which are reused across executions.

\textbf{Slot-grouped cuBLAS.} Packed blocks are grouped by identical
\((h_i,w_j)\) shape within each nonzero slot. One
\texttt{cublasSgemmGroupedBatched} call is issued per slot; outputs
within a call do not overlap, and later slots use \(\beta=1\) to
accumulate. Device pointer arrays for changing \(B\) and \(C\) addresses
are refreshed on the execution stream and this required marshaling is
included in timing.

The checked-in grid contains 128 workloads: 1,024 block rows and
columns; mean degrees \(\{2,4,8,16\}\); RHS widths \(\{8,16,32,64\}\);
four block-size distributions (uniform, low variance, high variance, and
bimodal); and local or random column selection. All runs use seed 1.
Each short-sweep point records seven measured iterations after three
warmups. GPU-event median and p95 times are reported separately from
synchronized host median and p95 times. Construction and host-to-device
input transfer are excluded.

The environment was an NVIDIA GeForce RTX 5060 Ti 16 GB, CUDA 13.4.46,
driver 610.57.01 with CUDA UMD 13.3, MSVC 19.44, Release configuration,
native architecture, and fast FP32 math. Profiler-instrumented durations
are not mixed with unprofiled benchmark timings.

Correctness is checked against a double-accumulating CPU reference using
a 64-case matrix spanning all RHS widths, distributions, degrees
\(\{1,4,8,16\}\), both locality patterns, and non-default streams. The
checked-in Release test passed on 21 August 2026. The project also
records successful memcheck and racecheck runs with zero findings.

\hypertarget{results}{%
\subsection{6. Results}\label{results}}

\hypertarget{end-to-end-regime-sweep}{%
\subsubsection{6.1 End-to-end regime
sweep}\label{end-to-end-regime-sweep}}

The final hybrid direct path is fastest among the four tested methods in
all 128 checked-in workloads. This win count is descriptive of this grid
and device, not evidence that the method dominates arbitrary matrices or
GPUs.

The optimization sequence is summarized below. Geometric means treat
every grid point equally and are not weighted to represent a production
workload distribution.

\begin{longtable}[]{@{}
  >{\raggedright\arraybackslash}p{(\columnwidth - 8\tabcolsep) * \real{0.18}}
  >{\raggedright\arraybackslash}p{(\columnwidth - 8\tabcolsep) * \real{0.18}}
  >{\raggedleft\arraybackslash}p{(\columnwidth - 8\tabcolsep) * \real{0.24}}
  >{\raggedleft\arraybackslash}p{(\columnwidth - 8\tabcolsep) * \real{0.24}}
  >{\raggedright\arraybackslash}p{(\columnwidth - 8\tabcolsep) * \real{0.18}}@{}}
\toprule
Change & Compared with & Evaluated slice & Geometric-mean median speedup
& Additional observation \\
\midrule
\endhead
Wider per-thread ILP & Four-accumulator direct kernel & All 128 cases &
1.24× & RHS 16/32/64: 1.21×/1.33×/1.43× \\
Cooperative \(B\) staging & Wider-ILP kernel & All 128 cases & 1.18× &
Every RHS-32 and RHS-64 case improved \\
Wider ILP + staging & Four-accumulator direct kernel & All 128 cases &
1.46× & Direct path retained 128/128 wins \\
RHS-32 double buffer + shape dispatch & Single-buffer staged kernel & 32
RHS-32 cases & 1.11× & All degree-8 and degree-16 cases improved \\
\bottomrule
\end{longtable}

Longer representative measurements provide scale beyond the short sweep:

\begin{longtable}[]{@{}
  >{\raggedright\arraybackslash}p{(\columnwidth - 6\tabcolsep) * \real{0.20}}
  >{\raggedleft\arraybackslash}p{(\columnwidth - 6\tabcolsep) * \real{0.27}}
  >{\raggedleft\arraybackslash}p{(\columnwidth - 6\tabcolsep) * \real{0.27}}
  >{\raggedleft\arraybackslash}p{(\columnwidth - 6\tabcolsep) * \real{0.27}}@{}}
\toprule
4,096-row workload & Final/accepted kernel & Comparison & Median
result \\
\midrule
\endhead
Bimodal, random, degree 16, RHS 64 & 2.895 ms staged & 4.887 ms
wider-ILP & 1.69× faster \\
Bimodal, random, degree 16, RHS 32 & 0.468 ms double-buffered & 0.549 ms
staged & 1.17× faster \\
\bottomrule
\end{longtable}

For the RHS-64 workload, the earlier four-accumulator result was 8.430
ms; the cumulative progression to 2.895 ms is 2.91×. These
representative figures answer an optimization question, but they should
not substitute for full long-duration reruns of every regime.

\hypertarget{profiling-evidence}{%
\subsubsection{6.2 Profiling evidence}\label{profiling-evidence}}

Profiling of the original difficult RHS-64/high-degree kernel showed
99.54\% achieved occupancy but only 0.48 eligible warps per scheduler
per cycle; 71.5\% of scheduler cycles had no eligible warp, and
long-scoreboard waits accounted for 90.1\% of the average issue
interval. DRAM throughput was 56.93\% of peak. Thus the limiting symptom
was load latency rather than insufficient occupancy, register spilling,
shared-memory spilling, or saturation of peak bandwidth.

Wider accumulation reduced redundant \(A\) traffic and raised useful
work per load. Cooperative staging then reduced the representative
long-scoreboard share from 78.0\% to 37.1\%, increased eligible warps
per scheduler from 0.41 to 0.75, and raised achieved occupancy from
49.6\% to 64.7\%. Barrier waiting became the dominant residual cost at
36.4\% of the issue interval.

For the final RHS-32 tall-row kernel, Nsight Compute confirms native
asynchronous global-to-shared loads. The replay-instrumented pass
reports 80.4\% achieved occupancy, 1.42 eligible warps per scheduler,
47.3\% of cycles with no eligible warp, 81.1\% SM throughput, and 68.1\%
DRAM throughput. These counters support the proposed mechanism, but
replay-instrumented time is excluded from benchmark speedups.

\hypertarget{discussion}{%
\subsection{7. Discussion}\label{discussion}}

The results illustrate three broader lessons. First, retaining block
structure can be valuable even when block dimensions are not
compile-time constants. Specializing a small orthogonal dimension---the
dense panel width---captures much of the benefit without a combinatorial
family of block-shape kernels. Second, occupancy alone did not diagnose
the slow kernel: high occupancy coexisted with few eligible warps and
dominant dependency stalls. Third, an optimization that is sound in
isolation need not be universal. Double buffering helped tall RHS-32
rows, harmed RHS 64 through shared-memory residency pressure, and was
not worth a second launch for low-degree mixed shapes.

The direct approach also has operational advantages. It owns outputs by
row, so it avoids floating-point atomics and their scheduling-dependent
accumulation order. It allocates no execute-time workspace, and
\texttt{execute} remains asynchronous on a caller-provided CUDA stream.
Conversely, plans rely on immutable matrix metadata, support only four
panel widths, and currently expose only FP32 non-transpose operation.
These choices make the implementation focused rather than
general-purpose.

\hypertarget{threats-to-validity-and-limitations}{%
\subsection{8. Threats to Validity and
Limitations}\label{threats-to-validity-and-limitations}}

The principal limitations are experimental.

\begin{itemize}
\tightlist
\item
  \textbf{One GPU and software stack.} Results may change with memory
  hierarchy, shared-memory capacity, scheduler behavior, library
  versions, or architecture.
\item
  \textbf{Synthetic matrices.} The generators control block-size
  distribution and column locality, but they do not establish
  performance on matrices from a real application.
\item
  \textbf{One seed and short timing runs.} Seven measured repetitions
  after three warmups are suitable for a development regime map, not
  publication-grade uncertainty estimates. The grid lacks confidence
  intervals, repeated process launches, and multi-seed aggregation.
\item
  \textbf{Restricted operation.} Only FP32, non-transpose \(A\),
  column-major dense panels, block dimensions 8--64 in steps of eight,
  and RHS widths 8--64 are supported.
\item
  \textbf{Baseline scope.} The baselines are credible persistent library
  compositions, but they are not an exhaustive comparison with research
  kernels, alternative reorderings, custom fixed-block ensembles, or
  application-specific approaches.
\item
  \textbf{Useful-GFLOP accounting.} Reported throughput counts
  arithmetic from stored dense blocks. It is useful for comparisons
  within this study but does not measure semantic sparsity inside a
  nominally dense stored block.
\item
  \textbf{Novelty scope.} Failure to identify a directly matching
  implementation in a scoped search is not proof that none exists.
\end{itemize}

Accordingly, the defensible conclusion is that VarBlockSpMM is a
well-supported single-GPU prototype whose design wins the tested regime
map---not that it is universally faster than cuSPARSE or cuBLAS.

\hypertarget{reproducibility}{%
\subsection{9. Reproducibility}\label{reproducibility}}

The repository requires CMake 3.25 or newer, CUDA 13.4, and an MSVC host
compiler with C++26 draft support. On Windows:

\begin{Shaded}
\begin{Highlighting}[]
\NormalTok{scripts\textbackslash{}build}\OperatorTok{.}\FunctionTok{ps1}
\NormalTok{build\textbackslash{}Release\textbackslash{}vbsr\_benchmark}\OperatorTok{.}\FunctionTok{exe} \OperatorTok{{-}{-}}\NormalTok{rows 4096 }\OperatorTok{{-}{-}}\NormalTok{degree 8 }\OperatorTok{{-}{-}}\NormalTok{rhs 32 }\OperatorTok{{-}{-}}\NormalTok{distribution high }\OperatorTok{{-}{-}}\NormalTok{locality random }\OperatorTok{{-}{-}}\NormalTok{warmup 10 }\OperatorTok{{-}{-}}\NormalTok{reps 50 }\OperatorTok{{-}{-}}\NormalTok{seed 1}
\NormalTok{scripts\textbackslash{}run\_grid}\OperatorTok{.}\FunctionTok{ps1} \OperatorTok{{-}}\NormalTok{Rows 4096 }\OperatorTok{{-}}\NormalTok{Reps 20 }\OperatorTok{{-}}\NormalTok{Warmup 5}
\end{Highlighting}
\end{Shaded}

The complete short sweeps are stored in \texttt{data/regime\_map.csv},
\texttt{data/regime\_map\_wide\_ilp.csv},
\texttt{data/regime\_map\_staged.csv}, and
\texttt{data/regime\_map\_double\_buffered.csv}. Nsight Compute exports
are retained separately under \texttt{data/}, and exact profiling
commands are documented in \texttt{docs/profiling.md}. A stronger
replication should run at least 20 measured iterations after five
warmups as already supported by the grid script, repeat each
configuration across independent process launches and seeds, report
bootstrap confidence intervals, and test multiple GPU architectures.

\hypertarget{conclusion}{%
\subsection{10. Conclusion}\label{conclusion}}

VarBlockSpMM demonstrates that packed variable-block SpMM can be mapped
efficiently to CUDA without global block padding, atomic output updates,
or execute-time temporary storage. The central mechanism is not a single
universal kernel but a restrained, measured policy: specialize
dense-panel width; reuse sparse values through independent accumulators;
stage dense-panel tiles only when reuse justifies synchronization; and
double-buffer only the row shapes and degrees that repay its resource
cost. On the tested RTX 5060 Ti grid, the final implementation beats its
scalar direct predecessor, persistent scalar-CSR cuSPARSE, and
slot-grouped cuBLAS in every case. The next scientific step is not
further headline optimization, but broader validation across hardware,
application matrices, seeds, and statistically rigorous runs.

\hypertarget{references}{%
\subsection{References}\label{references}}

{[}1{]} C. Yang, A. Buluç, and J. D. Owens, ``Design Principles for
Sparse Matrix Multiplication on the GPU,'' \emph{Euro-Par 2018}, 2018.
\url{https://arxiv.org/abs/1803.08601}

{[}2{]} G. Li et al., ``High Performance Unstructured SpMM Computation
Using Tensor Cores,'' 2024. \url{https://arxiv.org/abs/2408.11551}

{[}3{]} NVIDIA, ``cuSPARSE Generic API: \texttt{cusparseSpMM} and
Blocked-ELL,'' CUDA Toolkit Documentation.
\url{https://docs.nvidia.com/cuda/cusparse/}

{[}4{]} NVIDIA, ``cuBLAS
\texttt{cublas\textless{}t\textgreater{}gemmGroupedBatched},'' CUDA
Toolkit Documentation.
\url{https://docs.nvidia.com/cuda/cublas/\#cublas-t-gemmgroupedbatched}

{[}5{]} W. Boukaram, Y. Liu, P. Ghysels, and X. S. Li, ``Adaptive
Sketching Based Construction of H² Matrices on GPUs,'' \emph{IPDPSW
2025}, 2025. \url{https://arxiv.org/abs/2506.16759}

{[}6{]} NVIDIA, ``Asynchronous Data Copies,'' \emph{CUDA Programming
Guide}.
\url{https://docs.nvidia.com/cuda/cuda-programming-guide/04-special-topics/async-copies.html}

\end{document}
